%-----------------------------------------------------------
%   Capítulo 6 - Implementação
%-----------------------------------------------------------
\chapter{Implementação}
\label{cap:implementacao}

Este capítulo descreve como foram feitos o esquema, a placa, o \textit{firmware}
e a ferramenta do PC. Tudo está no repositório do projeto e é gerado ou
verificado por programas, para que cada alteração possa ser repetida e
conferida pela integração contínua (CI). A placa ainda não foi fabricada:
os ensaios em bancada ficam para o capítulo seguinte.

\section{Esquema elétrico}

O esquema tem sete folhas (alimentação, microcontrolador, Ethernet, CAN FD,
série, I2C e GPS, PWM), com 212 componentes. Cada folha é escrita por um
programa em Python (\texttt{hardware/gerador/folha\_*.py}) que coloca os
símbolos e liga os pinos por nome de rede. A atribuição de pinos do
microcontrolador fica numa única tabela, usada também para verificar o
\textit{firmware} (Secção~\ref{sec:impl-firmware}). A CI corre o verificador
de regras elétricas (ERC) do KiCad~9 em cada alteração; fica um único erro,
aceite, nos dois pinos VCAP unidos, como pede a ST. O esquema completo está
no Anexo~\ref{anexo:EETS}.

A lista de material sai do esquema em duas versões: uma com fabricantes e
preços para o orçamento (\texttt{hardware/bom/bom.csv}) e outra no formato da
montagem da JLCPCB, com o código de cada peça no catálogo LCSC. Ao preparar a
segunda, três peças estavam sem stock e foram substituídas por equivalentes
com o mesmo \textit{footprint} (a bobina de 2,2~µH, os cristais de 25~MHz e as
resistências de 60,4~$\Omega$ da terminação CAN, que passaram a 62~$\Omega$).
O isolador ADuM6401 e o conector DE-9 também não estão em stock e são
soldados à mão.

\section{Placa de circuito impresso}

A placa (100~$\times$~100~mm, quatro camadas) é gerada por
\texttt{hardware/gerador/pcb.py} a partir do esquema. As peças principais
têm posição fixa; os passivos são colocados junto do pino do circuito
integrado a que pertencem. O encaminhamento foi feito por uma sequência de
programas, cada um a partir do resultado do anterior
(Tabela~\ref{tab:impl-encaminhamento}).

\begin{table}[htbp]
\centering\small\setstretch{1.05}
\renewcommand{\arraystretch}{1.15}
\caption{Etapas do encaminhamento da placa.}
\label{tab:impl-encaminhamento}
\begin{tabularx}{\textwidth}{@{}>{\raggedright\arraybackslash}p{0.22\textwidth}>{\raggedright\arraybackslash}X>{\raggedleft\arraybackslash}p{0.13\textwidth}@{}}
\toprule
\rowcolor{LasetTint}
Programa & O que faz & Em falta \\
\midrule
\texttt{pcb.py}, \texttt{fanout.py} & disposição, zonas de potência e vias aos planos interiores & 302 \\
\texttt{pf.py} & encaminhamento por congestionamento negociado em F.Cu e B.Cu & 80 \\
\texttt{encaminhar.py} & pesquisa A* das ligações em conflito, também com pistas de 0,15~mm & 21 \\
\texttt{interior.py}, \texttt{melhorar.py} & sinais lentos pela camada de 3,3~V; arranque e novo encaminhamento das redes vizinhas & 11 \\
fecho por coordenadas & últimas ligações, com o afastamento conferido pelo DRC & 0 \\
\bottomrule
\end{tabularx}
\end{table}

Uma revisão independente da placa levou a mudanças no andar de potência: o
fusível passou a 125~V, o limiar de arranque baixou para 4,0~V, os MOSFET
passaram a ser de nível lógico e o andar foi espelhado para cada perna ficar
junto dos seus pinos no LM5177, o que reduziu o comando de porta mais longo
de 52 para 22~mm. Os nós de comutação e o relógio de 50~MHz do RMII ficaram
fora da camada interior. Depois disso:

\begin{itemize}
  \item \texttt{vias\_pad.py} tirou de dentro dos pads 78 das 90 vias que os
        encaminhadores lá tinham deixado (uma via aberta num pad puxa a
        solda); as 11 restantes ficam cheias e tapadas no fabrico;
  \item \texttt{baixo.py} passou cinco condensadores para a face de baixo, por
        baixo do pino: os do VCAP ficaram a 1,0 e 1,5~mm (estavam a 7 e 8~mm)
        e os do transceiver CAN~1 a 1,0 a 2,5~mm (estavam a 7 a 12~mm);
  \item \texttt{pilha.py} escreve na placa a pilha de camadas JLC04161H-7628,
        para a qual as larguras e afastamentos foram pensados.
\end{itemize}

A CI corre o verificador de regras de desenho (DRC) e a comparação com o
esquema no KiCad~9. Os ficheiros de fabrico (Gerber, furação, posições e
lista de material) saem de \texttt{fabrico.py} e \texttt{jlcpcb.py} e estão no
repositório, com os passos da encomenda em \texttt{ENCOMENDA.md}, na pasta
\texttt{hardware/fabrico}.

\section{\textit{Firmware}}
\label{sec:impl-firmware}

O \textit{firmware} (cerca de 2100 linhas de C) segue a arquitetura da
Secção~\ref{sec:proj-software}: as rotinas de receção colocam mensagens numa
fila de entrada, uma tarefa de encaminhamento aplica a tabela de regras e
cada saída tem uma fila urgente e uma normal. Usa FreeRTOS e LwIP, só com
memória estática, e compila para a NUCLEO-H753ZI e para a placa do projeto
com uma opção (\texttt{PLACA}). A Tabela~\ref{tab:impl-firmware} resume o
estado de cada parte.

\begin{table}[htbp]
\centering\small\setstretch{1.05}
\renewcommand{\arraystretch}{1.15}
\caption{Estado do \textit{firmware}.}
\label{tab:impl-firmware}
\begin{tabularx}{\textwidth}{@{}>{\raggedright\arraybackslash}X>{\raggedright\arraybackslash}p{0.34\textwidth}@{}}
\toprule
\rowcolor{LasetTint}
Parte & Estado \\
\midrule
Relógios a 480~MHz; FDCAN a 80~MHz, 1 + 8~Mbit/s com BRS & feito \\
Encaminhamento, filas urgente/normal, contadores de descarte & feito, com testes \\
Ligação ao PC por RS232 (COBS + CRC-16) & feito, com testes \\
Ligação ao PC por Ethernet (UDP e TCP) & esboço \\
Ligação ao PC por CAN FD (ISO 15765-2) & esboço, com testes \\
Configuração na \textit{flash} com dupla cópia (RF18), protegida contra cópias meio escritas & esboço, com testes \\
Vigilante independente (reinício automático em 4~s) & feito \\
RS485, I2C, PWM a 50~Hz com paragem por \textit{hardware} & feito \\
PPS no TIM2 e NMEA & feito, com testes \\
\bottomrule
\end{tabularx}
\end{table}

Os módulos que não dependem do \textit{hardware} (encaminhamento, política de
filas, COBS, CRC-16, NMEA, ISO-TP, configuração e formato das tramas) são
compilados também para o PC e testados na CI. Um teste lê uma trama gerada
pela ferramenta do PC, o que garante que os dois lados usam o mesmo formato.
Um teste de robustez passa milhões de entradas aleatórias ou corrompidas pelos
módulos que recebem dados de fora, compilado com verificação de acessos à memória
(AddressSanitizer): nenhuma escrita fora dos limites e nenhuma trama corrompida
aceite. Outro teste (\texttt{firmware/testes/pinos.py}) confere cada pino
configurado no \textit{firmware} com a tabela de pinos do esquema.

Ao rever o código foram corrigidos quatro casos que só se veriam na placa: o PWM
ficava desligado depois de uma falha transitória na entrada de paragem, o CAN FD
não recuperava de um \textit{bus-off}, o RS485 reencaminhava o eco do que a própria
unidade emitia, e uma configuração meio gravada (falta de alimentação durante a
gravação) bloqueava o arranque.

\section{Ferramenta do PC}

\texttt{firmware/ferramentas/laset\_pc.py} envia e recebe mensagens por RS232,
UDP, TCP ou CAN FD, com o formato da Tabela~\ref{tab:proj-mensagem}. Serve
para os ensaios (enviar uma trama CAN, ler um registo I2C, mudar a largura
de um impulso PWM) e para configurar a unidade: regras de encaminhamento,
ligação ao PC, endereço IP e velocidades das linhas série, gravadas na
\textit{flash} com \texttt{cfg-gravar}.
